%% kapa: exam=2 round=29 year=2018 period=2 subject=appraisal_theory subject_ko=감정평가이론 seq=2 number=2 points=30
%% source: Q-Net 감정평가사 기출문제 (cst00309) artlSeq=5203398 file=raw/1008/5203398/2교시 감정평가이론.pdf
%% conversion: kordoc 4.24.51 -> tools/gen2.py (auto); review=proofread-gpt-terra; flags=none
\begin{longquestion}{【문제 2】}
다음의 제시된 자료를 참고하여 물음에 답하시오. (30점) \par
\begin{framedtext}
인구 1,000만의 대도시인 A시와 약 40분 거리에 있는 인구 30만 규모의 기성도시인 B도시를 연결하는 전철이 개통되었다. 전철의 개통은 B도시의 광역접근성 개선효과를 가져와 부동산시장 및 부동산가격에 변화를 줄 것으로 예상된다.
\end{framedtext}
1) B도시에 새롭게 신설된 전철역세권의 지역분석에 대하여 설명하시오. (15점) \par
2) 전철개통으로 인한 접근성의 개선이 B도시의 유형별 부동산시장에 미치는 긍정적ㆍ부정적 효과에 대하여 설명하시오. (15점) \par
\end{longquestion}
